\documentclass[10pt]{beamer}

% =========================
% Thème Beamer
% =========================
\usetheme{Frankfurt}        % sobre et lisible
\usecolortheme{default}

% =========================
% Packages utiles
% =========================
\usepackage[french]{babel}
\usepackage[T1]{fontenc}
\usepackage{lmodern}


\RequirePackage{tabularray}
\RequirePackage{xcolor}

% Supprimer "Table X:" dans les légendes tabularray
\DefTblrTemplate{caption-tag}{default}{}
\DefTblrTemplate{caption-sep}{default}{}
\DefTblrTemplate{caption-text}{default}{}


\newcommand{\ok}[1]{\textcolor{teal}{#1}}
\newcommand{\nok}[1]{\textcolor{red}{#1}}
\newcommand\mydots{\makebox[1em][c]{.\hfil.\hfil.}}

% =========================
% Infos du document
% =========================
\title[Cas pratiques]{Cas pratiques juridiques: 1ère journée}
\author{Aurélien Gabieli}
\date{\today}

% =========================
% Footer : numéro de slide
% =========================
\setbeamertemplate{footline}[frame number]

% =========================
% Sommaire automatique
% =========================

% Slide plan au début
\begin{document}

\begin{frame}
  \titlepage
\end{frame}

\begin{frame}{Plan}
  \tableofcontents
\end{frame}

% =========================
% Plan au début de chaque section
% =========================
\AtBeginSection[]{
  \begin{frame}{Plan}
    \tableofcontents[currentsection]
  \end{frame}
}

% ======================
% Cas pratiques (1 à 12)
% ======================

% =========================
% Cas pratique 1
% =========================

\section[Cas 1]{Cas pratique 1 : Les constats d’huissier}

\begin{frame}{Question 1}

\begin{block}{}
L'art. \textbf{L615-5} al. 1 CPI : "tous les moyens".
\end{block}

\begin{table}[h]
\centering
\begin{tabular}{l l l}
    Quand? & Moyen de preuve \\
    \hline
    av. action en cf. & constat (- force probante: témoignage) \\
    au début d'une action en cf. & saisie cf. \\
    à la fin d'une action en cf. & droit de l'information \\
\end{tabular}
\end{table}
\end{frame}

\begin{frame}{Question 2}

\begin{block}{}
L'art. \textbf{D631-2} CPI : "Tribunal Judiciaire de Paris".
\\ Principe de droit: specialia generalibus derogant.
\end{block}

\begin{itemize}
    \item En l'espèce, Tribnual Jurdicaire de Saint-Dié-des-Vosges.
    De plus, L615-5 (spécifique) déroge à Art. 145 CPC (général).
    \item En conclusion, mauvaise compétence et fondement nul. 
    Donc, ordonnance nulle.
    Par conséquent, PV 14/11/2023 nul.
\end{itemize}
\end{frame}

\begin{frame}{Question 3}

\begin{itemize}
    \item En l'espèce, \begin{itemize}
        \item description et capture pages web pas ds pt. 1 ordonnance
        \item 1 PV ("je reprends"), pas plusieurs dans pt. 2 ordonnance.
        \item stagiaire, pas homme de l'art, pt. 3 ordonnance.
    \end{itemize}
    \item En conclusion, au moins pour ces raisons,
    le PV 15/11/2023 est nul.
\end{itemize}

\end{frame}

\begin{frame}{Question 4}

\begin{itemize}
    \item En l'espèce, analyse pas dans ordonnance.
    \item En conclusion, au moins pour cette raison,
    le PV 17/11/2023 est nul.
\end{itemize}

\end{frame}
% =========================
% Cas pratique 2
% =========================

\section[Cas 2]{Cas pratique 2 : La saisie-contrefaçon}

%--------------------------
\begin{frame}{Question 1}

\begin{block}{}
Art. \textbf{D211-6} COJ dispose que le tribunal judiciaire ayant compétence exclusive pour connaître des actions en matière de brevets d'invention, de certificats d'utilité, de certificats complémentaires de protection et de topographies de produits semi-conducteurs, dans les cas et conditions prévus par le code de la propriété intellectuelle, est celui de Paris.
\end{block}

\begin{itemize}
  \item En l'espèce, la saisie-contrefaçon en brevet est ordonnée par le Tribunal judiciaire de Paris.
  De plus, France Tube se prévaut d'une action en contrefaçon de la partie française d'un brevet européen.
  \item En conclusion, le tribunal jurdiciaire de Paris est compétent pour ordonner la saisie-contrefaçon.
\end{itemize}

\end{frame}
%--------------------------

%--------------------------
\begin{frame}{Question 2}

\begin{block}{}
Art. \textbf{L615-5} al. 3 CPI dispose que la juridiction peut ordonner, aux mêmes fins probatoires, la description détaillée ou la saisie réelle 
[...]
\end{block}

\begin{itemize}
  \item En l'espèce, la saisie intégrale du stock n'est pas une saisie probatoire.
  \item En conclusion, la saisie du stock ne semble pas licite.
\end{itemize}

\end{frame}

\begin{frame}{Question 2}
    
\begin{block}{}
Art. \textbf{496} CPC dispose que s'il est fait droit à la requête, tout intéressé peut en référer au juge qui a rendu l'ordonnance.
\end{block}
\begin{itemize}
    \item En l'espèce, le juge a fait droit à la requête de saisie-contrefaçon.
     Mais, la saisie ordonnée présente un élément illicite.
    \item En conclusion, la constestation se fait par la procédure de référé.
\end{itemize}
\begin{block}{}
Art. \textbf{497} CPC dispose que le juge a la faculté de modifier ou de rétracter son ordonnance, même si le juge du fond est saisi de l'affaire.
\end{block}
\begin{itemize}
    \item En l'espèce, la contestation semble légitime.
    Le juge a un motif pour rétracter l'ordonnance.
    \item En conclusion, le pv de la saisie est susceptible d'être annulé. 
\end{itemize}
    
\end{frame}
%--------------------------

%--------------------------
\begin{frame}{Question 3}
    
\begin{block}{}
\textbf{Principe de loyauté} (CA 06/04/2022 7ème page 2e par.): "ne donnant pas au juge des requêtes
tous les éléments en sa possession pour qu'il puisse exercer son contrôle".
\end{block}

\begin{itemize}
    \item En l'espèce, l'information sur l'opposition est pertinente pour ordonner la saisie-contrefaçon.
    Il y a en effet un sursit à statuer dans ce cas-là (Art. 378 CPC).
    \item En conclusion, l'ordonnance est susceptible d'être rétractée.
\end{itemize}

\end{frame}
%--------------------------

%--------------------------
\begin{frame}{Question 4}

\begin{block}{}
L'art. \textbf{60} AJUB dispose que la Juridication peut, avant même
l'engagement d'une action au fond, ordonner des mesures provisoires rapides
et efficaces pour conserver les éléments de preuve pertinents au regard
de la contrefaçon alléguée [...]
\end{block}

\begin{itemize}
  \item En l'espèce, France Tube a la possibilité de solliciter devant la JUB une mesure de préservation de preuve.
  \item En conclusion, France Tube a à sa disposition une mesure équivalente à la saisie-contrefaçon française.
\end{itemize}

\end{frame}

%--------------------------

\begin{frame}{Question 4}
  
\begin{block}{}
L'art. \textbf{60} AJUB dispose que à la demande du requérant
qui a présenté des éléments de preuve raisonnablement accessibles
pour étayer ses allégations selon lesquelles sont brevet a
été contrefait ou qu'une contrefaçon est imminente [...].
\end{block}

\begin{itemize}
  \item En l'espèce, il y a trois procès-verbal.
  \item En conclusion, une condition pour requérir à une mesure de préservation de preuve
  est remplie.
\end{itemize}

  \begin{table}[h]
\centering
\begin{tabular}{l l}
  Saisie CF JUB & Saisie CF FR \\
  \hline
  + portée territoriale & - portée territoriale\\
  procédure contradictoire & non contradictoire \\
  motifs requis (R192 RPJUB) & pas de preuve requise \\
\end{tabular}
\end{table}
\end{frame}
% =========================
% Cas pratique 3
% =========================

\section[Cas 3]{Cas pratique 3 : Confidentialité en
contrefaçon}

%--------------------------
\begin{frame}{Question 1: hypothèse rien de prévu}

La question propose le découpage suivant: 
Ensemble de documents "informations" = documents pas confidentiels $\cup$
documents confidentiels.\\
Objet: documents pas confidentiels.

\begin{block}{}
Interprétation stricte de l'ordonnance (\textbf{pourvoi n°03-12162}):
appartenait à l'huissier d'effectuer les opérations de saisie 
conformément aux dispositions de l'ordonnance le saisissant
\end{block}

\begin{itemize}
    \item En l'espèce,
    \begin{itemize}
        \item Hypothèse 1: saisie documents pas confidentiels $\in$ ordonnance.
        Non, pas de mention dans le sujet.
        \item Hypothèse 2: saisie documents pas confidentiels $\notin$ ordonnance.
        Oui, car rien de prévu (d'autre) dans l'ordonnance.
        Or, documents pas confidentiels $\neq$ échantillons de tube.
    \end{itemize}
    \item En conclusion, saisie illicite. 
    La saisie est licite si l'ordonnance vérifie l'hypothèse 1.
\end{itemize}

\end{frame}
%--------------------------

%--------------------------
\begin{frame}{Question 1: hypothèse séquestre provisoire}

Objet: documents pas confidentiels.

\begin{block}{}
Art. \textbf{R615-4} al. 4 CPI dispose que "Afin d'assurer la protection du secret des affaires, 
le président peut ordonner d'office le placement sous séquestre provisoire des pièces saisies, 
dans les conditions prévues à l'article R. 153-1 du code de commerce.".
\end{block}

\begin{itemize}
    \item En l'espèce, documents pas confidentiels.
    \begin{itemize}
        \item Hypothèse 1: saisie document pas confidentiels $\in$ séquestre provisoire. 
        \item Hypothèse 2: saisie document pas confidentiels $\notin$ séquestre provisoire.
    \end{itemize}
    Aucun élément ne permet de discrimer entre les deux hypothèses.
    \item En conclusion,
    \begin{itemize}
        \item Conclusion 1: la saisie est illicite.
        \item Conclusion 2: la saisie est licite.
    \end{itemize}
    La saisie est licite si le séquestre provisoire vérifie l'hypothèse 2.
\end{itemize}

\end{frame}
%--------------------------

%--------------------------
\begin{frame}{Question 2: hypothèse rien de prévu}

Ensemble de documents = documents pertinents $\cup$ documents pas pertinents.\\
Objet: ensemble de documents.

\begin{block}{}
Art. \textbf{L615-5} al. 2 3e phrase CPI dispose que "L'ordonnance peut autoriser 
[ci-après: première condition (i)] la saisie réelle de tout document se rapportant aux produits ou procédés prétendus contrefaisants
[ci-après: deuxième condition (ii)] en l'absence de ces derniers.".
\end{block}

\begin{itemize}
    \item On suppose que l'ordonannce a prévu ce cas-là (sinon 
    saisie illicite). Le stock de tube
        contrefaisant est 50\% de son chiffre d'affaire.
        Donc, ensemble de documents comprend probablement
        des documents pas pertinents. Donc,
        documents pas pertinents $\in$ saisie.
        Donc, première condition (i) pas vérifiée.
        De plus, il y a des stocks de tubes.
        Donc, deuxième condition (ii) pas vérifiée.

    \item En conclusion, la saisie est illicite.
    Pour s'oppposer à la saisie pendant la saisie-cf., 
    vérifier que l'ordonnance a prévu ce cas-là.
    Pour contester ensuite auprès du juge,
    vérifier la première condition (i)
    et vérifier la deuxième condition (ii).
\end{itemize}

\end{frame}

%--------------------------

%--------------------------

\begin{frame}{Question 2: hypothèse séquestre provisoire}


Objet: ensemble de documents.

\begin{block}{}
Art. \textbf{R615-4} al. 4 CPI dispose que "Afin d'assurer la protection du secret des affaires, 
le président peut ordonner d'office le placement sous séquestre provisoire des pièces saisies, 
dans les conditions prévues à l'article R. 153-1 du code de commerce.".
\end{block}

\begin{itemize}
    \item En l'espèce,
    ensemble de documents = documents pas confidentiels $\cup$ documents confidentiels.
    \item En conclusion, pas d'opposition pendant la saisie-cf.
    Toutefois, si documents pas confidentiels $\in$ séquestre provisoire,
    alors, l'ordonnance est nulle.\\    
    Donc, on comprend que le breveté n'a aucun intérêt à demander
    un séquestre provisoire dans l'ordonnance.
\end{itemize}

\end{frame}
%--------------------------

%--------------------------
\begin{frame}{Question 3}

\begin{block}{}
Art. \textbf{R153-1} al. 2 Code de Commerce: Si le juge n'est pas saisi d'une demande
 de modification ou de rétractation 
de son ordonnance en application de l'article 497 du code de procédure civile 
dans un délai d'un mois à compter de la signification de la décision, 
la mesure de séquestre provisoire mentionnée à l'alinéa précédent est levée 
et les pièces sont transmises au requérant.
\end{block}

\begin{itemize}
    \item Hypothèse rien de prévu: Asie Trading notifie le commissaire
    de justice que les documents sont confidentiels. Les documents
    sont placé au séquestre provisoire.
    \item Hypothèse séquestre provisoire: Les documents
    sont placés d'office au séquestre provisoire.
\end{itemize}
Dans les deux cas, la partie saisie doit saisir le juge dans un délai de 1 mois.
    Sinon, il y a levée du sequestre provisoire.
\end{frame}

%--------------------------
% =========================
% Cas pratique 4
% =========================

\section[Cas 4]{Cas pratique 4 : Recevabilité en
contrefaçon}

%--------------------------
\begin{frame}{Question 1}

\begin{longtblr}{colspec = {X[3,l] X[7,l] X[5,l]}, rowsep = 1pt}

\textbf{personne}
 & \textbf{condition?}
 & \\
\hline
\ok{déposant} et/ou \mydots
 & -- pas de condition
 & L615-2 al. 1 \\

\ok{co-pro} et/ou \mydots
 & -- notification aux copropriétaires
 & L613-29 \\

\ok{cessionnaire} et/ou \mydots
 & -- cession inscrite au RNB pour cessionnaire par transmission à titre particulier ou par transmission à titre universel
(p. ex: fusion, absorption, héritage)
 & L613-9 al. 1 \\

 & -- cession inscrite au RNB en procédure: régularisation rétroactive jusqu'au jour de la cession
 & Pourvoi n°D22-22.999 \\
\end{longtblr}

\end{frame}

\begin{frame}{Question 1}

\begin{longtblr}{colspec = {X[3,l] X[7,l] X[5,l]}, rowsep = 1pt}
\ok{licencié exclusif} et/ou \mydots
 & -- (i) notification au breveté et (ii) pas d'interdiction d'engager une action en cf.
 pour action en cessation de l'atteinte et action en réparation
 & L615-2 al. 2 \\

\ok{licencié non exclusif} et/ou \mydots
 & -- (i) notification au breveté et (ii) autorisation expression d'engager une action en contrefaçon
 %pour action en réparation de son préjudice ou intervention dans l'instance engagée par le breveté
 %pour réparation de son préjudice propre
 & L615-2 al. 6 \\

\ok{licencié}
 & -- pas de condition d'inscription au RNB pour se joindre à l'action du breveté
pour réparation de son préjudice propre
 & L613-9 al. 3\\
% &
% - pour les marques/modèles: licencié non inscrit au RNB engage une action &
% Cour de Justice 22/06/2016\\

\nok{cédant}
 & -- réparation du préjudice lorsque cession des droits de poursuite antérieurs
 & \\

\end{longtblr}

\end{frame}
%--------------------------

%--------------------------
\begin{frame}{Question 2}

\begin{block}{}
Art. \textbf{L615-2} al. 1 CPI: L'action en contrefaçon est exercée par le titulaire du brevet.
\end{block}

\begin{itemize}
    \item En l'espèce, brevet A déposé par la société France Tube.
    Pas de condition pour le déposant.
    \item En conclusion, France Tube est recevable à agir en contrefaçon du brevet A.
\end{itemize}

\begin{block}{}
Art. \textbf{L613-9} al. 1 CPI: Tous les actes transmettant ou modifiant 
les droits attachés à une demande de brevet ou à un brevet doivent, 
pour être opposables aux tiers, être inscrits sur un registre, 
dit Registre national des brevets, tenu par l'Institut national de la propriété industrielle.
\end{block}

\begin{itemize}
    \item En l'espèce, brevet B déposé par M. Gendre dirigeant de France Tube.
    \item En conclusion, France Tube est recevable à agir en contrefaçon du brevet A à
    condition d'inscrire la cession au RNB.
\end{itemize}

\end{frame}
%--------------------------

%--------------------------
\begin{frame}{Question 3}

\begin{block}{}
Art. \textbf{L615-2} al. 2 CPI: Sauf stipulation contraire du contrat de licence,
elle est également ouverte au titulaire d'une licence exclusive à condition,
à peine d'irrecevabilité, d'informer au préalable le titulaire du brevet.

\end{block}

\begin{itemize}
    \item En l'espèce, France Tube est licenciée exclusive du brevet C.
    \item En conclusion, France Tube est recevable à agir en contrefaçon du brevet C à
    condition de notifier le breveté et de ne pas avoir d'interdiction d'engager
    une action en contrefaçon.
\end{itemize}

\end{frame}
%--------------------------

%--------------------------
\begin{frame}{Question 4}

\begin{block}{}
Art. \textbf{47}(1) AJUB dispose que le titulaire du brevet européen a qualité pour agir
    en contrefaçon devant la Juridiction Unifiée du Brevet. \\
Art. \textbf{47}(2) AJUB Le licencié exclusif peut également agir, saus si l'accord
de licence en dispose autrement et à condition d'informer le titulaire du brevet.
\end{block}

\begin{itemize}
    \item En l'espèce, France Tube aurait souhaité agir devant la JUB car ses brevets EP
    n'ont pas fait l'objet d'une déclaration de dérogation (opt-out).
    \item En conclusion, devant la JUB, France Tube est recevable à agir en contrefaçon du brevet A et B.
    France Tube est recevable à agir en contrefaçon du brevet C à
    condition de notifier le breveté et de ne pas avoir d'interdiction d'engager
    une action en contrefaçon.
\end{itemize}

\begin{itemize}
    \item Les règles JUB sont proches de celles du droit français.
\end{itemize}

\end{frame}
%--------------------------
% =========================
% Cas pratique 5
% =========================

\section[Cas 5]{Cas pratique 5 : La copropriété du brevet}

%--------------------------
\begin{frame}{Question 1}

\begin{block}{}
L'article \textbf{L613-29}b) CPI dispose que 
chacun des copropriétaires peut agir en contrefaçon à son seul profit. 
Le copropriétaire qui agit en contrefaçon doit notifier l'assignation délivrée
 aux autres copropriétaires ; il est sursis à statuer sur 
l'action tant qu'il n'est pas justifié de cette notification.\\
+ \textbf{L615-2} (titulaire apte à agir en cf. )
+ \textbf{L613-32} (pas de stipulation contraire)
+ \textbf{L614-9} (traduction FR de demande EP).
\end{block}

\begin{itemize}
  \item En l'espèce, la société DUVIN est copropriétaire du brevet EP avec Mme RAISIN.
  La société DUVIN est apte à agir en cf. Il n'y pas de stipulation contractuelles
  contraires puisque pas de convention de copropriété.
  \item En conclusion, la société DUVIN a le droit d'agir seule
  devant le tribunal judiciaire de Paris, à condition de notifier Mme. RAISIN.
\end{itemize}

\end{frame}

%--------------------------
\begin{frame}{Question 2}

\begin{block}{}
L'article \textbf{L615-7} CPI dispose que les dommages-intérêts sont fixés par:\\
a) manque à gagner, perte subies;\\
b) préjudice moral;\\
c) bénéfices du cf.\\
+ \textbf{L615-3} CPI (tout moyens)
+ \textbf{L615-5-2} CPI (droit de l'information).
\end{block}

\begin{itemize}
  \item En conclusion, la société DUVIN a à sa disposition
  pour jusitifier des D\&I du droit de l'information.
\end{itemize}

\end{frame}

\begin{frame}{Question 2}

\begin{block}{}
L'article \textbf{1240} du Code Civil dispose que
tout fait quelconque de l'homme, qui cause à autrui un dommage,
oblige celui par la faute duquel il est arrivé, à le réparer.
\end{block}

\begin{itemize}
  \item En conclusion, la société Duvin ne peut pas demander
  la réparation du préjudice qui est propre à Mme. Raisin.
  \item Si Mme. Raisin souhaite obtenir réparation,
  Mme. Raison peut agir en réparation de son préjudice propre
  par \textbf{L613-29}b) CPI.
\end{itemize}

\end{frame}

%--------------------------
\begin{frame}{Question 3}

\begin{block}{}
L'article \textbf{L614-12} CPI renvoie à l'article 138 CBE.\\
+ article \textbf{70} CPC (lien suffisant)\\
+ \textbf{CA Paris, pôle 5 ch. 2, 13 avril 2012, n°11/01329}\\
+ \textbf{CA Lyon, 12 sept. 2013, n°11/06334}.\\
\end{block}

\begin{itemize}
    \item Par hypothèse, la société Aubonpain est assignée.
    \item En conclusion, la société Aubonpain a le droit de
    demander une demande reconventionnelle en nullité à condition
    de la diriger contre l'ensemble des titulaires du brevet,
    à savoir Mme. Raisin et la société Duvin. A defaut, la
    demande reconventionnelle est nulle.
\end{itemize}

\end{frame}

%--------------------------
\begin{frame}{Question 4}

\begin{block}{}
L'article \textbf{L613-29}c) CPI dispose que chacun des copropriétaires 
peut concéder à un tiers une licence d'exploitation non exclusive à son profit,
 sauf à indemniser équitablement les autres copropriétaires 
 qui n'exploitent pas personnellement l'invention ou 
 qui n'ont pas concédé de licence d'exploitation.
  A défaut d'accord amiable,
cette indemnité est fixée par le tribunal judiciaire.
\end{block}

\begin{block}{}
L'article \textbf{L613-29}d) CPI dispose que'une 
licence d'exploitation exclusive 
ne peut être accordée qu'
avec l'accord de tous les copropriétaires ou par autorisation de justice. 
\end{block}

\begin{itemize}
    \item En conclusion, la société Duvin a le droit de concéder à 
    la société Aubonpain
     une licence d'exploitation non exclusive à condition d'indemniser
    Mme. Raisin. Toutefois, la société Duvin n'a pas le droit de concéder
    à la société Aubonpain une licence exclusive sans l'accord de Mme. Raisin.
\end{itemize}

\end{frame}
% =========================
% Cas pratique 6
% =========================

\section[Cas 6]{Cas pratique 6 : Droit conférés par une demande de brevet}

%--------------------------
\begin{frame}{Question 1}

\begin{block}{}
L'article \textbf{L613-1} CPI dispose que le droit exclusif d'exploitation mentionné 
à l'article L. 611-1 prend effet à compter du dépôt de la demande.\\
+ \textbf{L615-4} (opposable à partir de la publication)
+ \textbf{L615-3} (interdiction provisoire) "toute personne habilitée
à agir en contrefaçon".
\end{block}

\begin{itemize}
  \item En l'espèce, la demande de brevet EP n'a pas été délivrée, mais
  la demande de brevet EP a été publiée.
  \item En conclusion, la société Hélicopathe est recevable pour agir
  en contrefaçon. La société hélicopathe a le droit de demander
  une interdiction provisoire.
\end{itemize}

\begin{block}{}
L'article \textbf{62} AJUB dispose que la Juridiction peut
prononcer des injections visant à 
à interdire, à titre provisoire et
sous réserve, le cas échéant, du paie
ment d'une astreinte, que la contrefaçon
présumée se poursuive.
\end{block}

\end{frame}

%--------------------------
\begin{frame}{Question 2}

\begin{itemize}
  \item En l'espèce, la demande de brevet EP
  la portée de la demande de brevet EP a été réduite.
  \item En conclusion, les chances délivrances d'un
  brevet EP sont plus grandes mais risque que le produit
  contrefaisant "sorte" de la portée du brevet EP délivré.
\end{itemize}

\end{frame}

%--------------------------
\begin{frame}{Question 3}

\begin{block}{}
L'article \textbf{L615-4} dispose que
par exception aux dispositions de l'article L. 613-1,
 les faits antérieurs à la date à laquelle
  la demande de brevet a été rendue publique 
  en vertu de l'article L. 612-21 ou
   à celle de la notification à tout tiers 
   d'une copie certifiée de cette demande 
   ne sont pas considérés comme 
ayant porté atteinte aux droits attachés au brevet. 
\end{block}

\begin{itemize}
  \item En l'espèce, la demande de brevet EP
  a été publiée.
  \item En conclusion, les actes de contrefaçon sont
  comptabilisés à partir de la date de publication
  de la demande de brevet EP.
\end{itemize}

\end{frame}

%--------------------------
\begin{frame}{Question 4}

\begin{block}{}
\begin{itemize}
    \item TGI de Paris: Article \textbf{378 CPC} (sursis pour bonne administration de justice).
    \item JUB: Article \textbf{33}(10) AJUB dispose que
les parties informent la Juridiction
de toute procédure de nullité, de limitation 
ou d'opposition pendante devant
l'Office européen des brevets, ainsi que
de toute demande de procédure accélérée
 présentée auprès de l'Office européen des brevets.
  La Juridiction peut
suspendre la procédure lorsqu'une
décision rapide peut être attendue de
l'Office européen des brevets.
\end{itemize}
\end{block}

\begin{itemize}
    \item En conclusion, devant TGI de Paris,
    un sursis à statuer est à prévoir.
    Devant la JUB, le sursis à statuer
    dépend de l'avancement de l'opposition.
\end{itemize}

\end{frame}


% =========================
% Cas pratique 7
% =========================

\section[Cas 7]{Cas pratique 7 : Procédure de limitation}

%--------------------------
\begin{frame}{Question 1}

\begin{block}{}
L'article \textbf{L613-24} al. 1 CPI dispose que
le propriétaire du brevet peut à tout moment soit renoncer 
à la totalité du brevet ou à une ou plusieurs revendications, 
soit limiter la portée du brevet 
en modifiant une ou plusieurs revendications.\\
L'article \textbf{L613-24} al. 3 CPI dispose que
le directeur de l'Institut national de 
la propriété industrielle examine 
la conformité de la requête avec
 les dispositions réglementaires 
 mentionnées à l'alinéa précédent.
\end{block}

\begin{itemize}
  \item En conclusion, l'INPI peut statuer
  sur la partie française du brevet EP.
  L'INPI peut statuer rapidement car il s'agit
  d'un examen formel.
\end{itemize}

\end{frame}

%--------------------------
\begin{frame}{Question 2}

\begin{block}{}
Cours de cassation 30/05/2012. 
\end{block}

\begin{itemize}
  \item En conclusion, le recours de Mycos
  contre la décision du Président de l'INPI
  est irrecevable. Néanmoins, Mycos a le droit
  d'engager une action en nullité contre
  la partie française du brevet EP.
\end{itemize}

\end{frame}

%--------------------------
\begin{frame}{Question 3}

\begin{block}{}
L'article \textbf{378} CPC (sursis pour bonne
adminisation de la justice).
\end{block}

\begin{itemize}
  \item Mycos a intérêt à demander un sursis à statuer.
  En effet, l'action en contrefaçon est sans fondement
  lorsque l'action en nullité = nullité.
  \item Phytos a intérêt à demander un sursis à statuer.
  En effet, il y a monopole d'exploitation tant que
  l'action en nullité est en cours.
\end{itemize}

\end{frame}


%--------------------------
\begin{frame}{Question 4}

\begin{block}{}
\textbf{Principe de proportionnalité} (CA n° 071/2022): Enfin, le principe de proportionnalité impose au juge de s'assurer qu'il existe un juste 
équilibre entre des droits fondamentaux opposés, en l'occurrence la loyauté des preuves 
dont dépend le respect du droit au procès équitable, d'une part, et le droit de propriété des 
titulaires de droits de propriété intellectuelle qui doit leur permettre de réunir des preuves, 
dans des conditions qui ne soient pas inutilement complexes ou coûteuses, afm d'assurer 
le respect de ces droits, d'autre part. 
\end{block}

Les critères sont:
\begin{itemize}
  \item le respect du droit au procès équitable;
  \item le droit de prioriété des titulaires de droits
  de propriété intellectuelle.
\end{itemize}

\end{frame}
% =========================
% Cas pratique 8
% =========================

\section[Cas 8]{Cas pratique 8 : Interdiction provisoire}

%--------------------------
\begin{frame}{Question 1}

\begin{block}{}
L'art. \textbf{L615-3} CPI dispose que "peut saisir en référé
la juridication civile compétente" ou "la juridication civile
compétente peut également ordonner [...] sur requête".
\end{block}

\begin{itemize}
  \item En l'espèce, Phyto souhaite l'une quelconque des
  interdictions provisoire.
  \item En conclusion, Phyto a la possibilité de s'adresse
  au juge des référés ou au juge de mise en l'état.   
\end{itemize}

\end{frame}
%--------------------------


%--------------------------
\begin{frame}{Question 2}

\begin{block}{}
L'art. \textbf{L615-3} al. 1 3e phrase CPI dispose que
"preuve raisonnablement accessible au demandeur et
vraisemblable qu'il est porté atteinte à ses droits"
\end{block}

\begin{itemize}
  \item En l'espèce, Phyto estime qu'il est porté
  atteinte à certaines revendications de son brevet.
  On suppose qu'une preuve raisonnablement accessible
  au demandeur est disponible. Il n'y a pas d'urgence
  apparente.
  \item En conclusion, la procédure en référé est
  préférée.
\end{itemize}

\end{frame}
%--------------------------

%--------------------------
\begin{frame}{Question 3}

\begin{block}{}
L'art. \textbf{L615-3} al. 1 3e phrase CPI ne dispose pas
d'une condition sur le statut du brevet.
\end{block}

\begin{itemize}
  \item En l'espèce, le brevet est l'objet d'un recours.
  \item En conclusion, une mesure provisoire n'est pas
  interdite.
\end{itemize}

\end{frame}
%--------------------------

%--------------------------
\begin{frame}{Question 4}

\begin{block}{}
\textbf{Principe de proportionnalité} (Ordonnance JME TGI Paris, 19 avril 2013, page 5) :
"raisonnablement pas ignorer la nullité [des revendications]". 
\end{block}

\begin{itemize}
  \item En l'espèce, le brevet est maintenu après la 1ère instance sous forme limitée.
  La validité du brevet est vraisemblable.
  \item En conclusion, les chances de succès d'une demande d'interdiction provisoire sont bonnes.
\end{itemize}

\end{frame}
%--------------------------

%--------------------------
\begin{frame}{Question 5}

\begin{block}{}
Art. 490 CPC (référé) ou Art. 496 CPC (requête) + A539 CPC (principe)
\\ Exception (Ordonnance TGI Paris, 15 avril 2016): "indépendamment de
l'absence de caractère suspensif de ce recours".
\end{block}

\begin{itemize}
  \item Par hypothèse, une interdiction provisoire est ordonnée.
  Hypothèse 1: en référé; Hypothèse 2: sur requête.
  \item En conclusion, dans les 2 cas, un recours est possible,
  mais pas suspensif.
\end{itemize}

\end{frame}
%--------------------------
% =========================
% Cas pratique 9
% =========================

\section[Cas 9]{Cas pratique 9 : Droit de possession
personnelle antérieure}

%--------------------------
\begin{frame}{Question 1}

\begin{block}{}
L'article \textbf{L613-3}(a) CPI (actes de contrefaçon)
+ \textbf{L615-1 al. 3} CPI (conditions de responsabilité)
+ \textbf{L615-4} CPI (opposable à partir de la publication).
\end{block}

\begin{itemize}
  \item En l'espèce, Padrole a un brevet \textbf{français} delivré.
  \begin{itemize}
    \item L'élément matériel: la revendication de produit sur la
  pince.
    \item L'élément légal: acte de fabrication de Mercurey,
  acte d'offre, de mise sur le marché et détention en vue
  de l'offre et de la mise sur le marché de Langsamer.
  Padrole bénéficie de la connaissance de cause, mais
  pas Mercurey.
    \item L'élément temporel: publication de la demande de brevet
  en août 2015. Les actes de contrefaçon ne sont comptabilisés
  qu'à partir de ce mois-là. Or, la commercialisation a
  commencé en janvier 2015. 
  \end{itemize}
  \item En conclusion, un risque de contrefaçon pour Padrole
  et Mercurey existe sur le territoire français.
\end{itemize}

\end{frame}

%--------------------------
\begin{frame}{Question 2}

\begin{block}{}
L'article \textbf{L611-11 CPI} (nouveauté)
+ \textbf{L611-13 CPI} (état de la technique).
\end{block}

\begin{itemize}
  \item En l'espèce, les prototypes dès décembre 2013 < dépôt FR.
  Mais une clause de confidentialité.
  \item En conclusion, raisonnement sur défaut de nouveauté
  pas de fondement. A moins d'une diffusion illicite
  prouvée avant la date de dépôt.
\end{itemize}

\end{frame}

%--------------------------
\begin{frame}{Question 3}

\begin{block}{}
L'article \textbf{L613-7 CPI} (droit
de possession personnelle antérieure) dispose
d'une exception à L613-3 CPI.
\begin{itemize}
  \item une possession de l’invention sur le territoire français,
  \item une exploitation effective ou des préparatifs sérieux,
  \item une possession de bonne foi.
\end{itemize}
\end{block}  

\begin{itemize}
  \item En l'espèce, \begin{itemize}
    \item Mercurey possède l'invention, pas Langsamer.
    \item 3 ans de R\&D.
    \item Pas d'élément de mauvaise foi.
  \end{itemize}
  \item En conclusion, Mercurey, pas Langsamer, a le droit de se
  prévaloir d'une possession personnelle antérieure.
\end{itemize}

\end{frame}
% =========================
% Cas pratique 10
% =========================

\section[Cas 10]{Cas pratique 10 : Contrefaçon
par fourniture de moyens}

%--------------------------
\begin{frame}{Question}

\begin{block}{}
L'article \textbf{26} AJBU (droit d'empêcher d'exploitation
indirecte de l'invention).
\end{block}

\begin{longtblr}{colspec = {X[1,l] X[3,l] X[3,l]}, rowsep = 0pt}

Critère & Contrefaçon & Pas contrefaçon \\

\hline
territoralité
& Papgeneric livre et offre de livrer sur
le territoire des Etats membres contractants
(implicite)
& \\

moyens essentiels
& du point de vue juridique oui, la revendication 1
comprend cette caractéristique
& du point de vue
technique non, l'invention fonctionne sans papier
avec des coins spéciaux dans un mode dégradé \\

personne habilité
& une personne autre que celle habilitée
à exploiter l'invention brevetée (implicite).
& \\

moyens aptes et destinés
& moyens aptes oui 
& moyens destinés difficile à vérifier. \\

\end{longtblr}

\end{frame}
% =========================
% Cas pratique 11
% =========================

\section[Cas 11]{Cas pratique 11 : Médiation judiciaire}

%--------------------------
\begin{frame}{Question 1}

\begin{longtblr}{colspec = {X[1,l] X[2,l] X[2,l]}, rowsep = 3pt}

Art. \textbf{750-1} CPC
& \textbf{Avantages}
& \textbf{Inconvénients} \\

\hline
\textbf{Conciliation} (art. \textbf{21} CPC)
&
Très rapide \newline
Faible coût \newline
Adaptée aux litiges simples
&
Peu efficace en contentieux technique complexe \newline
Accord souvent difficile à obtenir
\\

\textbf{Médiation judiciaire} (art. \textbf{131-1} et s.)
&
Confidentialité \newline
Souplesse des échanges \newline
Utile pour négocier licence/transaction
&
Coût du médiateur \newline
Dépend de la coopération des parties
\\

\textbf{Procédure participative}
(art. \textbf{2062} s. C.civ. ; art. \textbf{1542} s. CPC)
&
Encadrée par les avocats \newline
Travail technique structuré \newline
Peut accélérer la résolution
&
Nécessite forte collaboration \newline
Peu utilisée en pratique
\\

\end{longtblr}

\end{frame}

%--------------------------
\begin{frame}{Question 2}

\begin{block}{}
Loi n°95-125 du 8 février 1995 (médiation judiciaire).
+ Art. \textbf{131-1} à \textbf{131-15} CPC.
+ art. \textbf{131-14} CPC et art. \textbf{21-3} CPC
(confidentialité)
+ art. \textbf{131-9} CPC (information au juge).
\end{block}

\begin{itemize}
  \item \textbf{Entretiens séparés :}
  le médiateur peut organiser des caucus pour recevoir des informations
  qu’une partie ne souhaite pas révéler immédiatement.

  \item \textbf{Secret et loyauté :}
  le médiateur est tenu à la confidentialité et les parties doivent participer
  de bonne foi.

  \item \textbf{Risques si échec :}
  les échanges restent en principe inutilisables au procès,
  mais l’information peut être découverte par d’autres voies
  (expertise, communication de pièces), avec un risque stratégique pour FUNITALIA.
\end{itemize}

\end{frame}

%--------------------------
\begin{frame}{Question 3}

\begin{longtblr}{colspec = {X[1,l] X[1,l]}, rowsep = 1pt}

\textbf{DUTCHJETSKI} & \textbf{FUNITALIA} \\

\hline

Réduire l’aléa judiciaire (nullité + débat technique complexe). 
& Éviter une interdiction et des dommages-intérêts élevés. \\

Obtenir une solution globale (France/Italie + ventes Espagne). 
& Régler confidentiellement les ventes sous autre appellation. \\

Limiter les coûts et la durée (experts, saisies, contentieux long). 
& Négocier une licence, un calendrier ou un écoulement de stock. \\

Accéder à des informations commerciales via accord amiable. 
& Réduire les risques financiers et réputationnels. \\

\end{longtblr}

\end{frame}
% =========================
% Cas pratique 12
% =========================

\section[Cas 12]{Cas pratique 12 : Standard Essential Patent et
abus de position dominante}

%--------------------------
%--------------------------
\begin{frame}{Question}

\begin{block}{}
\begin{itemize}
  \item \textbf{Art. 102 TFUE} (abus de position dominante).
  \item \textbf{CJUE, 16 juillet 2015, Huawei Technologies / ZTE} (SEP + FRAND).
\end{itemize}
\end{block}

Se comporter comme un \textbf{willing licensee} (candidat sérieux à une licence FRAND),
afin de rendre une demande d'injonction \textbf{risquée / abusive} pour le titulaire du SEP
(standard essential patent).

\end{frame}

%--------------------------
\begin{frame}{Question}

Trois conditions d'un "wiling licensee":
\begin{itemize}
  \item \textbf{Réagir sans délai} dès la notification/ mise en demeure : manifester
  une volonté claire d'obtenir une licence \textbf{FRAND} (pas de stratégie dilatoire).
  
  \item \textbf{Négocier de bonne foi} : demander une offre FRAND, échanger sur la base
  (portefeuille, périmètre, taux, assiette, pays, durée), coopérer activement.
  
  \item \textbf{Si désaccord} : formuler une \textbf{contre-offre écrite} à des conditions FRAND
  et fournir une \textbf{garantie} (séquestre / garantie bancaire) pour les redevances,
  avec \textbf{reddition de comptes} des actes d'exploitation.
\end{itemize}

L'interdiction est évitable si TELEPHONICA
met en oeuvre les trois mesures préventives.

\end{frame}

\end{document}